
\documentclass{article}
\usepackage{colortbl}
\usepackage{makecell}
\usepackage{multirow}
\usepackage{supertabular}

\begin{document}

\newcounter{utterance}

\centering \large Interaction Transcript for game `cladder', experiment `full\_v1.5\_default', episode 9763 with Qwen3.5{-}4B.
\vspace{24pt}

{ \footnotesize  \setcounter{utterance}{1}
\setlength{\tabcolsep}{0pt}
\begin{supertabular}{c@{$\;$}|p{.15\linewidth}@{}p{.15\linewidth}p{.15\linewidth}p{.15\linewidth}p{.15\linewidth}p{.15\linewidth}}
    \# & \multicolumn{2}{c}{Player} && \multicolumn{2}{c}{Game Master} \\
    \hline

    \theutterance \stepcounter{utterance}  
    & & & \multicolumn{4}{p{0.6\linewidth}}{
        \cellcolor[rgb]{0.9,0.9,0.9}{
            \makecell[{{p{\linewidth}}}]{
                \texttt{\tiny{[P1$\langle$GM]}}
                \texttt{You are an expert in causal inference. The following question is not a typical commonsense query, but rather a meticulously designed question created by a professor specializing in causal inference, intended to assess the students' mastery of the course content.} \\
\texttt{Imagine a self{-}contained, hypothetical world with only the following conditions, and without any unmentioned factors or causal relationships: Yupt has a direct effect on jyka and kwox. Jyka has a direct effect on kwox. We know that yupt causes jyka. yupt and jyka causes kwox. We observed an individual is yupt. Would an individual is kwox if not jyka instead of jyka?} \\
\texttt{Start your answer with 'Yes' or 'No', followed by additional reasoning or evidence to support your explanation.} \\
            }
        }
    }
     \\ \\

    \theutterance \stepcounter{utterance}  
    & & \multicolumn{4}{p{0.6\linewidth}}{
        \cellcolor[rgb]{0.95,0.95,0.95}{
            \makecell[{{p{\linewidth}}}]{
                \texttt{\tiny{[P1$\leftrightarrow$P]}}
                \texttt{{[}DRAFT{]}} \\
\texttt{Yes} \\
\\ 
\texttt{**Reasoning:**} \\
\\ 
\texttt{To determine the answer, we must analyze the causal structure defined in the hypothetical world and apply the rules of causal inference regarding counterfactuals.} \\
\\ 
\texttt{1.  **Causal Structure**:} \\
\texttt{    *   \$Yupt \textbackslash{}rightarrow Jyka\$} \\
\texttt{    *   \$Yupt \textbackslash{}rightarrow Kwox\$} \\
\texttt{    *   \$Jyka \textbackslash{}rightarrow Kwox\$} \\
\texttt{    *   This forms a "fork" or "collider" structure depending on the direction of interest, but specifically, \$Kwox\$ is a common effect of both \$Yupt\$ and \$Jyka\$.} \\
\\ 
\texttt{2.  **The Observation**:} \\
\texttt{    *   We observed an individual where \$Yupt = 1\$ (True).} \\
\\ 
\texttt{3.  **The Counterfactual Question**:} \\
\texttt{    *   The question asks: "Would an individual be \$Kwox\$ if not \$Jyka\$ instead of \$Jyka\$?"} \\
\texttt{    *   Formally, this asks for the counterfactual outcome \$Kwox\_\{\textbackslash{}neg Jyka\}\$ given the observed world \$Yupt=1\$. In other words, if we intervene to set \$Jyka = 0\$ (while keeping \$Yupt=1\$), what is the value of \$Kwox\$?} \\
\\ 
\texttt{4.  **Causal Mechanism Analysis**:} \\
\texttt{    *   In this world, \$Kwox\$ is caused by two factors: \$Yupt\$ and \$Jyka\$.} \\
\texttt{    *   Since we observed \$Yupt=1\$, the causal pathway \$Yupt \textbackslash{}rightarrow Kwox\$ is active. Therefore, \$Yupt\$ contributes to making \$Kwox\$ true.} \\
\texttt{    *   The question posits a scenario where \$Jyka\$ is set to false (\$\textbackslash{}neg Jyka\$).} \\
\texttt{    *   Does removing the effect of \$Jyka\$ (\$Jyka \textbackslash{}rightarrow Kwox\$) negate the effect of \$Yupt\$ (\$Yupt \textbackslash{}rightarrow Kwox\$)?} \\
\texttt{    *   **No.** Causal effects are additive in this context (or at least, the presence of one cause does not eliminate the necessity of another for the effect to occur, unless specified otherwise as a "necessary and sufficient" condition, which is not stated). The prompt states "\$Yupt\$ and \$Jyka\$ causes \$Kwox\$," implying a conjunction of causes or a cumulative effect.} \\
\texttt{    *   Even if we interpret the relationship strictly logically: If \$A\$ causes \$C\$ and \$B\$ causes \$C\$, then \$A\$ alone is sufficient to cause \$C\$ (assuming no other conditions are required for \$A \textbackslash{}rightarrow C\$).} \\
\texttt{    *   Therefore, if \$Yupt\$ is present, \$Kwox\$ will occur regardless of whether \$Jyka\$ is present or absent.} \\
\\ 
\texttt{5.  **Conclusion**:} \\
            }
        }
    }
    & \\ \\

    \theutterance \stepcounter{utterance}  
    & & \multicolumn{4}{p{0.6\linewidth}}{
        \cellcolor[rgb]{0.95,0.95,0.95}{
            \makecell[{{p{\linewidth}}}]{
                \texttt{\tiny{[P1$\leftrightarrow$P]}}
                \texttt{{[}SELF{-}FEEDBACK{]}} \\
\texttt{FORMAT: problem} \\
\texttt{RULES: problem} \\
\texttt{STRATEGY: problem} \\
\texttt{REASON: The answer fails to follow the required format of starting with 'Yes' or 'No' followed by reasoning, instead providing a full paragraph with numbered lists and LaTeX formatting.} \\
            }
        }
    }
    & \\ \\

    \theutterance \stepcounter{utterance}  
    & \multicolumn{4}{p{0.6\linewidth}}{
        \cellcolor[rgb]{0.561,0.737,0.561}{
            \makecell[{{p{\linewidth}}}]{
                \texttt{\tiny{[P1$\rangle$GM]}}
                \texttt{Yes} \\
\texttt{Reasoning:} \\
\texttt{To determine the answer, we must analyze the causal structure defined in the hypothetical world and apply the rules of causal inference regarding counterfactuals.} \\
\\ 
\texttt{1. Causal Structure:} \\
\texttt{Yupt has a direct effect on jyka and kwox.} \\
\texttt{Jyka has a direct effect on kwox.} \\
\texttt{We know that yupt causes jyka.} \\
\texttt{yupt and jyka causes kwox.} \\
\\ 
\texttt{2. The Observation:} \\
\texttt{We observed an individual is yupt.} \\
\\ 
\texttt{3. The Counterfactual Question:} \\
\texttt{The question asks: Would an individual be kwox if not jyka instead of jyka?} \\
\texttt{Formally, this asks for the counterfactual outcome kwox if not jyka given the observed world yupt. In other words, if we intervene to set jyka = 0 (while keeping yupt = 1), what is the value of kwox?} \\
\\ 
\texttt{4. Causal Mechanism Analysis:} \\
\texttt{In this world, kwox is caused by two factors: yupt and jyka.} \\
\texttt{Since we observed yupt, the causal pathway yupt {-}> kwox is active. Therefore, yupt contributes to making kwox true.} \\
\texttt{The question posits a scenario where jyka is set to false (not jyka).} \\
\texttt{Does removing the effect of jyka negate the effect of yupt?} \\
\texttt{No. The prompt states "yupt and jyka causes kwox," implying that the presence of yupt is sufficient to cause kwox in conjunction with the other potential cause. Even if jyka is removed, the direct effect of yupt on kwox remains active.} \\
\texttt{Therefore, if yupt is present, kwox will occur regardless of whether jyka is present or absent.} \\
\\ 
\texttt{5. Conclusion:} \\
\texttt{Yes, an individual would be kwox if not jyka instead of jyka.} \\
            }
        }
    }
    & & \\ \\

    \theutterance \stepcounter{utterance}  
    & & & \multicolumn{2}{p{0.3\linewidth}}{
        \cellcolor[rgb]{0.9,0.9,0.9}{
            \makecell[{{p{\linewidth}}}]{
                \texttt{\tiny{[GM$|$GM]}}
                \texttt{Yes} \\
            }
        }
    }
    & & \\ \\

    \theutterance \stepcounter{utterance}  
    & & & \multicolumn{2}{p{0.3\linewidth}}{
        \cellcolor[rgb]{0.9,0.9,0.9}{
            \makecell[{{p{\linewidth}}}]{
                \texttt{\tiny{[GM$|$GM]}}
                \texttt{no} \\
            }
        }
    }
    & & \\ \\

    \theutterance \stepcounter{utterance}  
    & & & \multicolumn{2}{p{0.3\linewidth}}{
        \cellcolor[rgb]{0.9,0.9,0.9}{
            \makecell[{{p{\linewidth}}}]{
                \texttt{\tiny{[GM$|$GM]}}
                \texttt{game\_result = LOSE} \\
            }
        }
    }
    & & \\ \\

\end{supertabular}
}

\end{document}
